\documentclass{article}
\usepackage{amsmath, amssymb, amsthm, algorithm, algpseudocode, graphicx, bm, booktabs, xcolor}
\usepackage[margin=1in]{geometry}

% Custom commands
\newcommand{\E}{\mathbb{E}}
\newcommand{\R}{\mathbb{R}}
\newcommand{\norm}[1]{\|#1\|}
\newcommand{\inner}[2]{\langle #1, #2 \rangle}
\newcommand{\grad}{\nabla}
\newcommand{\ind}{\mathbb{I}}
\newcommand{\step}[1]{\text{[Step #1]}}
\newcommand{\vect}[1]{\bm{#1}}
\newcommand{\diag}{\text{diag}}
\newcommand{\argmin}{\text{argmin}}
\newcommand{\argmax}{\text{argmax}}
\newcommand{\Var}{\text{Var}}
\newcommand{\Cov}{\text{Cov}}
\newcommand{\tr}{\text{tr}}
\newcommand{\sign}{\text{sign}}

% Theorem environments
\newtheorem{assumption}{Assumption}
\newtheorem{theorem}{Theorem}
\newtheorem{lemma}{Lemma}
\newtheorem{corollary}{Corollary}
\newtheorem{definition}{Definition}
\newtheorem{remark}{Remark}
\newtheorem{proposition}{Proposition}

% Algorithm formatting
\algnewcommand{\LineComment}{\State \(\triangleright\) #1}
\algnewcommand{\LineCommentEnd}{\(\triangleright\) #1}

\begin{document}

\title{RMSProp: Algorithm Description, Convergence Theory, and Complete Proof for Non-Convex Smooth Functions}
\author{}
\date{}
\maketitle

% ============================================================
% 1. ALGORITHM DETAILS EXPANSION
% ============================================================

\section*{Algorithm Name}
\textbf{RMSProp (Root Mean Square Propagation)}

\section*{Mathematical Formulation}

Let \(f: \R^d \to \R\) be the objective function. At iteration \(t = 1,2,\dots,T\), the algorithm maintains a parameter vector \(\vect{w}_t \in \R^d\) and a running average of squared gradients \(\vect{v}_t \in \R^d\). Given a stochastic gradient oracle that returns \(\vect{g}_t = \grad f(\vect{w}_{t-1}; \vect{z}_t)\) where \(\vect{z}_t\) is a random sample, the updates are:

\begin{align}
\vect{v}_t &= \beta \vect{v}_{t-1} + (1-\beta) \vect{g}_t \odot \vect{g}_t, \label{eq:v_update} \\
\vect{w}_t &= \vect{w}_{t-1} - \alpha_t \frac{\vect{g}_t}{\sqrt{\vect{v}_t + \epsilon}}, \label{eq:w_update}
\end{align}

where:
\begin{itemize}
    \item \(\odot\) denotes element-wise (Had